\documentclass[10pt,a4paper]{amsart}
\usepackage{amsmath,amssymb,amsthm,graphicx}
\usepackage{verbatim}
\newenvironment{reference}{\comment}{\endcomment}
%\newenvironment{reference}{}{}
\newenvironment{slogan}{\comment}{\endcomment}
\newenvironment{history}{\comment}{\endcomment}

% For commutative diagrams we use Xy-pic
\usepackage[all]{xy}

% We use 2cell for 2-commutative diagrams.
\xyoption{2cell}
% Theorem environments.
%
\theoremstyle{plain}
\newtheorem{theorem}[subsection]{Theorem}
\newtheorem{proposition}[subsection]{Proposition}
\newtheorem{lemma}[subsection]{Lemma}

\theoremstyle{definition}
\newtheorem{definition}[subsection]{Definition}
\newtheorem{example}[subsection]{Example}
\newtheorem{exercise}[subsection]{Exercise}
\newtheorem{situation}[subsection]{Situation}

\theoremstyle{remark}
\newtheorem{remark}[subsection]{Remark}
\newtheorem{remarks}[subsection]{Remarks}

\numberwithin{equation}{subsection}

% Macros
%
\def\lim{\mathop{\mathrm{lim}}\nolimits}
\def\colim{\mathop{\mathrm{colim}}\nolimits}
\def\Spec{\mathop{\mathrm{Spec}}}
\def\Hom{\mathop{\mathrm{Hom}}\nolimits}
\def\Ext{\mathop{\mathrm{Ext}}\nolimits}
\def\SheafHom{\mathop{\mathcal{H}\!\mathit{om}}\nolimits}
\def\SheafExt{\mathop{\mathcal{E}\!\mathit{xt}}\nolimits}
\def\Sch{\mathit{Sch}}
\def\Mor{\mathop{\mathrm{Mor}}\nolimits}
\def\Ob{\mathop{\mathrm{Ob}}\nolimits}
\def\Sh{\mathop{\mathit{Sh}}\nolimits}
\def\NL{\mathop{N\!L}\nolimits}
\def\CH{\mathop{\mathrm{CH}}\nolimits}
\def\proetale{{pro\text{-}\acute{e}tale}}
\def\etale{{\acute{e}tale}}
\def\QCoh{\mathit{QCoh}}
\def\Ker{\mathop{\mathrm{Ker}}}
\def\Im{\mathop{\mathrm{Im}}}
\def\Coker{\mathop{\mathrm{Coker}}}
\def\Coim{\mathop{\mathrm{Coim}}}

% Boxtimes
%
\DeclareMathSymbol{\boxtimes}{\mathbin}{AMSa}{"02}

%
% Macros for moduli stacks/spaces
%
\def\QCohstack{\mathcal{QC}\!\mathit{oh}}
\def\Cohstack{\mathcal{C}\!\mathit{oh}}
\def\Spacesstack{\mathcal{S}\!\mathit{paces}}
\def\Quotfunctor{\mathrm{Quot}}
\def\Hilbfunctor{\mathrm{Hilb}}
\def\Curvesstack{\mathcal{C}\!\mathit{urves}}
\def\Polarizedstack{\mathcal{P}\!\mathit{olarized}}
\def\Complexesstack{\mathcal{C}\!\mathit{omplexes}}
% \Pic is the operator that assigns to X its picard group, usage \Pic(X)
% \Picardstack_{X/B} denotes the Picard stack of X over B
% \Picardfunctor_{X/B} denotes the Picard functor of X over B
\def\Pic{\mathop{\mathrm{Pic}}\nolimits}
\def\Picardstack{\mathcal{P}\!\mathit{ic}}
\def\Picardfunctor{\mathrm{Pic}}
\def\Deformationcategory{\mathcal{D}\!\mathit{ef}}
\usepackage{cleveref}
\crefname{originallemma}{original lemma}{original lemmas}
\Crefname{originallemma}{Original lemma}{Original lemmas}
\makeatletter
\newlabel{lemma-when-surjective-to-completion}{{090S}{source}{}{sourcepointer-1}{}}
\newlabel{lemma-when-surjective-to-completion@cref}{{[originallemma][0][]090S}{[1][0][]source}{}{}{}}
\makeatother
\title{Completeness for a larger ideal passes to any finitely generated subideal}
\author{The Stacks Project Authors}
\date{stacks-project:algebra.tex:local-snapshot-20261012; GNU FDL 1.2 or later}
\begin{document}
\maketitle
\noindent\textit{Editorial note.} All delivered mathematical text is the author's own, copied from the pinned Stacks Project chapter file algebra.tex; no statement or proof step was rewritten, reordered or removed.\par\medskip
\noindent\textit{Editorial note.} Preamble provenance, as required for this source: the frozen snapshot stages the chapter's own preamble file beside it (preamble.tex, 217 lines, 8116 bytes, the same shared Stacks-wide preamble the published chow and etale-cohomology cards use); the first line of the pinned chapter file is \textbackslash{}input\{preamble\}, so no preamble text lives inside the chapter itself. The card takes that file's lines 8-18 (the verbatim package and the three \textbackslash{}comment-based environments reference/slogan/history at lines 9-12, section 13.2, together with the Xy-pic setup) and its lines 154 to the runtime end of the file (the native theorem declarations and mathematical macros). Not one character was changed, no package was added, no macro was defined, and the excerpt serves only as preamble, never as a span coordinate.\par\medskip
\noindent\textit{Editorial note.} The same-chapter pointers are located at their \textbackslash{}label lines inside the pinned chapter file; the cross-chapter pointers are located in the project's own tag list staged in the snapshot (sources/stacks-chow/tags.txt), at the verbatim lines carrying their canonical tags.\par\medskip
\noindent\textit{Editorial note.} The counted proof refers to a same-chapter criterion by label, which this lane has itself claimed under its own card; it is recorded as a navigation stub located at its \textbackslash{}label line, and the argument consuming it stays inside the counted span.\par\medskip
\noindent\textit{Editorial note.} The counted spans contain no \textbackslash{}cite, so the card carries no bibliography.\par\medskip
\noindent\textit{Allowed standard prerequisites (editorial).}\begin{itemize}
\item Surjectivity onto the completion tested on one generator. Invoked in the counted proof at line 23002. Reduces the claim to the case where the smaller ideal is principal.
\item Completeness as convergence of series. Used in the counted proof at lines 23014-23018. Lets the geometric series in the chosen generator be summed to exhibit the required limit.
\end{itemize}\medskip
\ifdefined\hypertarget\hypertarget{sourcepointer-1}{}\fi
\noindent\textit{Source reference.} Original lemma 090S: sources/stacks-algebra/pinned-original-source.tex, source lines 22973--22973. Navigation pointer for \textbackslash{}ref\{lemma-when-surjective-to-completion\} used inside the counted span (line 23002); the argument that consumes it is retained in the span.\par\medskip
\begin{lemma}
\label{lemma-complete-by-sub}
Let $A$ be a ring. Let $I \subset J \subset A$ be ideals.
If $M$ is $J$-adically complete and $I$ is finitely generated, then
$M$ is $I$-adically complete.
\end{lemma}

\begin{proof}
Assume $M$ is $J$-adically complete and $I$ is finitely generated.
We have $\bigcap I^nM = 0$ because $\bigcap J^nM = 0$. By
Lemma \ref{lemma-when-surjective-to-completion}
it suffices to prove the surjectivity of $M \to \lim M/I^nM$ in case
$I$ is generated by a single element. Say $I = (f)$.
Let $x_n \in M$ with $x_{n + 1} - x_n \in f^nM$. We have to show there exists
an $x \in M$ such that $x_n - x \in f^nM$ for all $n$.
As $x_{n + 1} - x_n \in J^nM$ and as $M$ is $J$-adically complete,
there exists an element $x \in M$ such that $x_n - x \in J^nM$.
Replacing $x_n$ by $x_n - x$ we may assume that $x_n \in J^nM$.
To finish the proof we will show that this implies $x_n \in I^nM$.
Namely, write $x_n - x_{n + 1} = f^nz_n$.
Then
$$
x_n = f^n(z_n + fz_{n + 1} + f^2z_{n + 2} + \ldots)
$$
The sum $z_n + fz_{n + 1} + f^2z_{n + 2} + \ldots$ converges in $M$
as $f^c \in J^c$. The sum $f^n(z_n + fz_{n + 1} + f^2z_{n + 2} + \ldots)$
converges in $M$ to $x_n$ because
the partial sums equal $x_n - x_{n + c}$ and $x_{n + c} \in J^{n + c}M$.
\end{proof}

\end{document}
